\subsection{Subcategory projection and exact twist support}
\label{sec:subcategory-projection}

Information about the center of a proper fusion subcategory imposes constraints
on individual restriction columns of the ambient center. It does not require
a braiding on the ambient fusion category, or an extension of a chosen
subcategory braiding.

\begin{proposition}\label{prop:central-subcategory-projection}
Let $\mathcal D\subset\mathcal C$ be a full fusion subcategory of a pivotal
fusion category, with the inherited pivotal structure. Let $P(V)$ be the sum
of all simple summands of $V$ belonging to $\mathcal D$.
If $(V,c)$ is an object of $Z(\mathcal C)$, then $P(V)$ has a canonical
$\mathcal D$-half-braiding. If $V$ is central simple with twist $\lambda$,
every simple constituent of $P(V)$ in $Z(\mathcal D)$ has twist $\lambda$.
Consequently, writing $F_{\mathcal D}:Z(\mathcal D)\to\mathcal D$ for the
forgetful functor, the $\mathcal D$-part of a central restriction column
belongs to the nonnegative integral span of the columns
$[F_{\mathcal D}(W)]$ with $\theta_W=\lambda$.
\end{proposition}
\begin{proof}
The complement of $\mathcal D$ is stable under tensoring on either side by
objects of $\mathcal D$. Indeed, if a simple $S\notin\mathcal D$ and
$X,T\in\mathcal D$ satisfy $\operatorname{Hom}(T,X\otimes S)\ne0$,
rigidity makes $S$ a constituent of $X^*\otimes T\in\mathcal D$, a
contradiction. Hence $P$ is a projection of $\mathcal D$-bimodule categories.
Put $W=P(V)$ and choose $i:W\to V$, $p:V\to W$ with $pi=1_W$.
For $X\in\mathcal D$ the maps
\[
 c^W_X=(1_X\otimes p)c_X(i\otimes1_X)
\]
are isomorphisms. Naturality and the half-braiding identity restrict to this
summand, so they make $W$ an object of $Z(\mathcal D)$.

It remains to check the twist, rather than merely the half-braiding.
Suppress associators and use the Drinfeld isomorphism
\[
 u_V=(\operatorname{ev}_V\otimes1_{V^{**}})
       (c_{V^*}\otimes1_{V^{**}})
       (1_V\otimes\operatorname{coev}_{V^*}):V\longrightarrow V^{**}.
\]
The duality identities give
\[
 \begin{split}
 (1_{V^*}\otimes p^{**})\operatorname{coev}_{V^*}
   &=(p^*\otimes1_{W^{**}})\operatorname{coev}_{W^*},\\
 \operatorname{ev}_V(p^*\otimes1_V)
   &=\operatorname{ev}_W(1_{W^*}\otimes p).
 \end{split}
\]
Using naturality of $c$ in its second argument between these two identities
gives $p^{**}u_Vi=u_W$. If $j$ denotes the pivotal isomorphism, then
$p^{**}j_Vi=j_W$. With the convention $\theta=u^{-1}j$, central simplicity
gives $u_V=\lambda^{-1}j_V$, and therefore
$u_W=\lambda^{-1}j_W$. Thus $\theta_W=\lambda1_W$.
Decomposing $W$ into central simples proves the integral-span assertion.
The opposite twist convention gives the same conclusion when used in both
centers.
\end{proof}

The projection also gives exact, twist-resolved dimension identities.
Here and below spherical structures are inherited, so the centers carry their
standard ribbon structures.

\begin{corollary}\label{cor:subcategory-twist-masses}
Let $\mathcal C$ be spherical and $\mathcal D\subset\mathcal C$ a full fusion
subcategory. Put $D_{\mathcal C}=\dim(\mathcal C)$ and
$D_{\mathcal D}=\dim(\mathcal D)$. For a simple $X\in\mathcal D$ and a root
of unity $\lambda$,
\[
 \sum_{\substack{Z\in\operatorname{Irr}Z(\mathcal C)\\\theta_Z=\lambda}}
       \dim(Z)[F_{\mathcal C}(Z):X]
 =\frac{D_{\mathcal C}}{D_{\mathcal D}}
   \sum_{\substack{W\in\operatorname{Irr}Z(\mathcal D)\\\theta_W=\lambda}}
       \dim(W)[F_{\mathcal D}(W):X].
\]
An empty right-hand side is zero. In the pseudo-unitary case all weights are
positive and dimensions can be replaced by Frobenius--Perron dimensions.
\end{corollary}
\begin{proof}
Choose $N$ divisible by the orders of all twists in $Z(\mathcal D)$.
Proposition~\ref{prop:central-subcategory-projection} shows that every central
simple contributing to the left-hand side has twist in $\mu_N$.
The indicator formula \cite{NS} gives, for $1\le n<N$,
\[
 D_{\mathcal C}\nu_n(X)
 =\sum_Z\theta_Z^n\dim(Z)[F_{\mathcal C}(Z):X].
\]
The same formula holds in $\mathcal D$, and the indicators agree: their
defining pivotal rotations act on the same spaces
$\operatorname{Hom}(1,X^{\otimes n})$. For $n=0$, the corresponding equality
is $\dim I_{\mathcal C}(X)=D_{\mathcal C}\dim(X)$, and similarly in
$\mathcal D$. Fourier inversion on $\mu_N$ proves the assertion.
\end{proof}

\begin{corollary}\label{cor:fibonacci-ising-projection}
In the pseudo-unitary setting the following subcategory constraints hold.
\begin{enumerate}
\item For a Fibonacci subcategory with basis $(1,a)$, $a^2=1+a$, write
$\varphi=\operatorname{FPdim}(a)=(1+\sqrt5)/2$. A nonzero projected
column $(r,s)$ has either twist $1$ and $r\ge s$, or $r=0$ and twist
$\zeta_5^2$ or $\zeta_5^{-2}$. For each of the latter twists, the
$a$-mass in Corollary~\ref{cor:subcategory-twist-masses} is
$D_{\mathcal C}\varphi/(\varphi+2)$; the twist-one $a$-mass is
$D_{\mathcal C}\varphi^2/(\varphi+2)$.
\item For an Ising subcategory with basis $(1,g,\sigma)$, a nonzero projected
column $(r,s,t)$ has one of the following forms: twist $1$ with $t=0$ and
$r\ge s$; twist $-1$ with $r=t=0$; or one of four Ising twists
$\{z,-z,z^{-1},-z^{-1}\}$ with $r=s=0$, where $z$ is an appropriate
primitive sixteenth root. The $g$-masses at $1$ and $-1$ are each
$D_{\mathcal C}/2$. The $\sigma$-mass at each of its four twists is
$D_{\mathcal C}\sqrt2/4$.
\end{enumerate}
In particular, a high column, whose unit coordinate is zero, cannot combine
$g$ and $\sigma$ in its Ising projection, or contain either $a$, $g$ or
$\sigma$ at twist one in the respective cases.
\end{corollary}
\begin{proof}
The rank-two and Ising classifications \cite{Ostrik2,EGNO} provide
nondegenerate braidings on these subcategories. For either such category
$\mathcal D$, its center is
$\mathcal D\boxtimes\mathcal D^{\mathrm{rev}}$ and the restriction of
$(U,V)$ is $U\otimes V$.
For Fibonacci the four columns are
\[
 (1,0),\quad(0,1),\quad(0,1),\quad(1,1),
\]
with twists $1,z,z^{-1},1$. Positivity gives the usual Fibonacci modular
$S$-matrix; balancing gives $z^2+\varphi z+1=0$, hence
$z=\zeta_5^{\pm2}$. The center dimensions are respectively
$1,\varphi,\varphi,\varphi^2$.
For Ising, the twist-one columns are $(1,0,0)$ twice and $(1,1,0)$ once;
the twist-$-1$ columns are $(0,1,0)$ twice; and each of
$z,-z,z^{-1},-z^{-1}$ has the column $(0,0,1)$ once. Their dimensions are
$1,2,\sqrt2$ as determined by these columns. Taking their additive spans
and applying the preceding corollary gives every assertion.
All inherited Ising spherical choices must be retained. If the second
indicator of $\sigma$ is fixed, the identity
$z^2+z^{-2}=\nu_2(\sigma)\sqrt2$ selects the appropriate set of four twists.
\end{proof}

\paragraph{Computational use.}
The criteria first restrict the allowed high-column/twist pairs and the
twist-one low columns. They then add the exact mass identities to the
integer linear equations. When using normalized relative twist traces,
one projects the identities to the chosen field, combining coincident
trace types; one must not confuse a projected trace with an actual twist.
These are necessary conditions only. General exclusions still require the
separately checked general-to-pseudo-unitary reduction.
